\documentclass[12pt,a4paper]{article}

\usepackage[margin=1in]{geometry}
\usepackage{graphicx}
\usepackage{amsmath,amssymb}
\usepackage{booktabs}
\usepackage{enumitem}
\usepackage{hyperref}
\usepackage{xcolor}
\usepackage{float}
\usepackage{caption}
\usepackage{subcaption}
\usepackage{algorithm}
\usepackage{algpseudocode}
\usepackage{listings}
\usepackage{multirow}
\usepackage{longtable}
\usepackage{fancyhdr}
\usepackage{setspace}

\hypersetup{
    colorlinks=true,
    linkcolor=blue!70!black,
    citecolor=green!50!black,
    urlcolor=blue!70!black,
}

\lstset{
    basicstyle=\ttfamily\small,
    breaklines=true,
    frame=single,
    backgroundcolor=\color{gray!10},
    keywordstyle=\color{blue!70!black},
    commentstyle=\color{green!50!black},
    stringstyle=\color{red!60!black},
}

\pagestyle{fancy}
\fancyhf{}
\fancyhead[L]{\small MineOptima --- Automated Underground Mine Design}
\fancyhead[R]{\small\thepage}
\renewcommand{\headrulewidth}{0.4pt}

\onehalfspacing

% ============================================================
\begin{document}

\begin{titlepage}
\centering
\vspace*{2cm}

{\Huge\bfseries MineOptima}\\[0.5cm]
{\LARGE Automated Underground Mine Design\\Using Operations Research, Machine Learning,\\and Monte Carlo Simulation}

\vspace{2cm}

{\Large\bfseries Technical Report}

\vspace{1.5cm}

{\large
\textbf{Author:} Mobashshir Zainuddin\\[0.2cm]
Roll No: 23MI3PE01\\[0.3cm]
Department of Mining Engineering\\
Indian Institute of Technology (IIT) Kharagpur\\[1cm]
\textbf{Date:} November 2026
}

\vspace{2cm}

\begin{abstract}
\noindent
This report presents \textbf{MineOptima}, an integrated computational platform for automated underground gold mine design. The system combines Binary Integer Linear Programming (BILP) for stope layout optimization, Random Forest machine learning for grade prediction, Monte Carlo simulation for risk analysis, production scheduling with Discounted Cash Flow (DCF) NPV, and geotechnical stability assessment using the Mathews method. The platform models an Archean lode gold deposit with geologically realistic log-normal grade distribution (CV $\sim$180\%), nugget effect, and vein-hosted structures---calibrated against published parameters from Canadian Shield gold operations. Built with a Flask + Python backend and Plotly.js frontend, the system replaces weeks of manual mine design iteration with a pipeline that runs in under 60 seconds. This report details the methodology, implementation, results, and future research directions.
\end{abstract}

\end{titlepage}

\tableofcontents
\newpage

% ============================================================
\section{Introduction}
\label{sec:introduction}

Underground mine design is a complex, multi-disciplinary engineering challenge that integrates geology, rock mechanics, economics, and operations research. Traditional practice relies on sequential manual workflows across disconnected commercial software packages (Surpac, Deswik, Datamine), where mine planners manually define stope boundaries through trial-and-error iteration over days to weeks.

This project develops \textbf{MineOptima}, an automated platform that integrates the entire underground mine design pipeline into a single web-based application. The key contributions are:

\begin{enumerate}[itemsep=0.3em]
    \item A geologically calibrated synthetic block model for Archean lode gold deposits with realistic nugget effect and vein-hosted grade structures.
    \item A floating stope algorithm coupled with Binary Integer Linear Programming (BILP) for provably optimal stope layout selection.
    \item Monte Carlo simulation for NPV risk quantification under price, grade, and cost uncertainty.
    \item Random Forest machine learning for grade prediction with per-block uncertainty quantification.
    \item Grade-priority production scheduling with Discounted Cash Flow analysis.
    \item Geotechnical stability assessment using the Mathews stability graph method and pillar Factor of Safety analysis.
    \item A responsive web dashboard with 9 interactive tabs for real-time parameter exploration.
\end{enumerate}

\subsection{Why Gold is Different}

Gold deposits---particularly orogenic (lode) gold---are fundamentally different from disseminated base-metal deposits such as copper, iron, or nickel. Gold occurs as free nuggets and fine particles within narrow quartz veins and shear zones, creating several unique modeling challenges:

\begin{itemize}[itemsep=0.2em]
    \item \textbf{Nugget effect:} Coarse visible gold creates extreme grade variability. The coefficient of variation (CV) for gold grades typically ranges from 100--300\%, compared to 40--80\% for copper \cite{journel1978}.
    \item \textbf{Log-normal distribution:} Gold grades follow a highly positively skewed log-normal distribution, with a few extremely high-grade samples dominating the metal content.
    \item \textbf{Narrow structures:} Ore zones are often narrow (1--5\,m true width), tabular veins with sharp ore/waste contacts, unlike the large disseminated bodies of porphyry copper deposits.
    \item \textbf{High-grade shoots:} Gold concentration occurs in plunging ore shoots at structural intersections (fold hinges, fault junctions), creating discrete targets within the vein system.
\end{itemize}

These characteristics directly influence block model design, stope optimization constraints, and economic evaluation---all of which are addressed in this platform.

% ============================================================
\section{Methodology}
\label{sec:methodology}

\subsection{Block Model Generation}
\label{sec:block-model}

The synthetic block model represents a typical Archean lode gold deposit hosted in a greenstone belt environment (analogous to Canadian Shield deposits such as Red Lake and LaRonde). The model domain consists of $40 \times 30 \times 20 = 24{,}000$ blocks, each measuring $5\,\text{m} \times 5\,\text{m} \times 5\,\text{m}$, covering a volume of $200\,\text{m} \times 150\,\text{m} \times 100\,\text{m}$ starting at 50\,m depth.

\subsubsection{Geological Structure}

The deposit model includes three structural elements:

\begin{enumerate}
    \item \textbf{Main quartz vein:} A narrow tabular body centered in the domain, with across-strike extent $\sigma_x = 15$\,m, along-strike $\sigma_y = 50$\,m, and down-dip $\sigma_z = 30$\,m. Peak grade 5.0\,g/t Au.
    \item \textbf{Hangingwall splay vein:} A secondary structure offset from the main vein ($\sigma = [10, 30, 18]$\,m), peak grade 3.5\,g/t Au.
    \item \textbf{High-grade shoot:} A plunging ore shoot at the intersection of the two structures ($\sigma = [12, 15, 10]$\,m), peak grade 8.0\,g/t Au---representing the structural trap where gold concentrates.
\end{enumerate}

\subsubsection{Grade Distribution}

The grade model captures gold's distinctive statistical properties:

\begin{equation}
    g_{\text{base}}(\mathbf{x}) = \sum_{k=1}^{3} A_k \exp\left(-\frac{1}{2} \sum_{i=1}^{3} \left(\frac{x_i - c_{k,i}}{\sigma_{k,i}}\right)^2 \right)
\end{equation}

where $A_k$ is the peak grade amplitude, $\mathbf{c}_k$ is the center, and $\boldsymbol{\sigma}_k$ controls the spatial extent of each structure. A spatially correlated noise field is generated via Gaussian filtering ($\sigma_{\text{filter}} = 1.5$ blocks) of random normal deviates, then added to the base grade.

The log-normal transformation is applied to capture gold's characteristic positive skew:

\begin{equation}
    g(\mathbf{x}) = \exp\left(\ln(g_{\text{base}}(\mathbf{x}) + 0.01) + \epsilon \right), \quad \epsilon \sim \mathcal{N}(0, 0.4)
\end{equation}

\subsubsection{Nugget Effect}

The nugget effect---representing coarse visible gold---is modeled by applying random exponential bonuses to approximately 2\% of blocks within vein structures:

\begin{equation}
    g_{\text{nugget}} \sim \text{Exponential}(\lambda = 8.0 \text{ g/t}) \quad \text{with probability } p = 0.02 \text{ (in-vein blocks)}
\end{equation}

This produces occasional extreme values (up to $\sim$40\,g/t), consistent with the erratic assay results characteristic of gold deposits.

\subsubsection{Rock Types and Alteration}

Blocks are classified into four rock types reflecting the alteration zonation around gold-bearing structures:

\begin{itemize}[itemsep=0.1em]
    \item \textbf{Quartz Vein} --- ore-bearing vein material ($\rho = 2.65$\,t/m$^3$)
    \item \textbf{Altered Basalt} --- proximal alteration envelope
    \item \textbf{Sericite Zone} --- intermediate alteration halo
    \item \textbf{Fresh Basalt} --- unaltered host rock ($\rho = 2.80$\,t/m$^3$)
\end{itemize}

\subsubsection{Calibration Parameters}

\begin{table}[H]
\centering
\caption{Block model parameters and calibration sources}
\begin{tabular}{@{}lll@{}}
\toprule
\textbf{Parameter} & \textbf{Value} & \textbf{Reference} \\
\midrule
Block size & $5 \times 5 \times 5$\,m & Industry standard for UG stopes \\
Grade CV & $\sim$180\% & Gold typical: 100--300\% \\
Grade max & $\sim$39\,g/t & Nugget effect (coarse gold) \\
Median grade & $\sim$0.39\,g/t & Log-normal positive skew \\
Mean grade & $\sim$0.94\,g/t & Archean lode gold deposits \\
Density (vein) & 2.65\,t/m$^3$ & Quartz vein material \\
Density (host) & 2.80\,t/m$^3$ & Mafic volcanic host rock \\
RQD & 65 $\pm$ 8\% & Precambrian hard rock \\
\bottomrule
\end{tabular}
\label{tab:calibration}
\end{table}

% -----------------------------------------------------------
\subsection{Economic Evaluation}
\label{sec:economics}

Each block is assigned an economic value based on contained metal and costs:

\begin{equation}
    V_j = T_j \times \left( g_j \times R \times P - C_m - C_p - C_{\text{ga}} \right)
    \label{eq:block-value}
\end{equation}

where $T_j$ is the tonnage (block volume $\times$ density), $g_j$ is gold grade (oz/t, converted from g/t by dividing by 31.1035), $R = 0.92$ is metallurgical recovery, $P = \$1{,}900$/oz is gold price, $C_m = \$45$/t mining cost, $C_p = \$25$/t processing cost, and $C_{\text{ga}} = \$10$/t general and administrative cost.

The breakeven cutoff grade is:
\begin{equation}
    g_{\text{cutoff}} = \frac{C_m + C_p + C_{\text{ga}}}{R \times P} \times 31.1035 \approx 1.42 \text{ g/t}
\end{equation}

Blocks with $g_j \geq g_{\text{cutoff}}$ are classified as \textit{ore}; the remainder as \textit{waste}.

% -----------------------------------------------------------
\subsection{Stope Optimization}
\label{sec:stope-opt}

\subsubsection{Candidate Generation: Floating Stope Algorithm}

Candidate stopes are generated by sweeping a rectangular template across the block model domain. For each position, the algorithm evaluates all feasible stope geometries within user-defined constraints:

\begin{itemize}[itemsep=0.1em]
    \item Width: 2--3 blocks (10--15\,m)
    \item Height: 3--5 blocks (15--25\,m)
    \item Length: 2--4 blocks (10--20\,m)
\end{itemize}

Each candidate stope records its enclosed blocks, total economic value, tonnage, metal content, average grade, and ore fraction.

\subsubsection{BILP Formulation}

The stope layout optimization is formulated as a Binary Integer Linear Program:

\begin{align}
    \text{Maximize:} \quad & Z = \sum_{i \in \mathcal{S}} V_i \cdot x_i \label{eq:bilp-obj} \\
    \text{Subject to:} \quad & \sum_{i \in \mathcal{S}(j)} x_i \leq 1 \quad \forall\, j \in \mathcal{B} \label{eq:bilp-overlap} \\
    & x_i \in \{0, 1\} \quad \forall\, i \in \mathcal{S} \label{eq:bilp-binary}
\end{align}

where $\mathcal{S}$ is the set of candidate stopes, $V_i$ is the total economic value of stope $i$, $x_i$ is the binary decision variable (1 = selected, 0 = not selected), $\mathcal{B}$ is the set of all blocks, and $\mathcal{S}(j)$ is the set of candidate stopes containing block $j$.

Constraint~\eqref{eq:bilp-overlap} ensures no two selected stopes share a block (non-overlap / pillar constraint). The problem is solved using \textbf{PuLP 4.x} with the \textbf{HiGHS} open-source mixed-integer linear programming solver, providing provably optimal solutions within a configurable time limit.

% -----------------------------------------------------------
\subsection{Mine Layout Generation}
\label{sec:mine-layout}

The development network is algorithmically generated from the optimized stope layout:

\begin{enumerate}
    \item \textbf{Decline:} A spiral ramp from surface to the deepest mining level (decline gradient $1:7$, development cost \$8{,}000/m).
    \item \textbf{Level drives:} Horizontal tunnels at each mining level connecting all stopes on that level (\$5{,}000/m).
    \item \textbf{Crosscuts:} Short perpendicular tunnels connecting level drives to individual stopes (\$5{,}000/m).
\end{enumerate}

Total development cost is computed as the sum of all segment lengths multiplied by their per-metre costs.

% -----------------------------------------------------------
\subsection{Monte Carlo NPV Simulation}
\label{sec:monte-carlo}

The Monte Carlo module quantifies project risk by simulating NPV under correlated input uncertainty:

\begin{equation}
    \text{NPV}_k = \sum_{i \in \mathcal{S}^*} T_i \times \left( g_i \cdot \tilde{f}_g \times R \times \tilde{P}_k - \tilde{C}_k \right) - C_{\text{dev}}
\end{equation}

where $\tilde{P}_k \sim \mathcal{N}(P, \sigma_P^2)$, $\tilde{f}_g \sim \mathcal{N}(1, \sigma_g^2)$ is a grade scaling factor, and $\tilde{C}_k$ includes cost variability. The simulation runs 1{,}000+ iterations and reports:

\begin{itemize}[itemsep=0.1em]
    \item Mean and standard deviation of NPV
    \item Probability of positive NPV: $P(\text{NPV} > 0)$
    \item Percentiles (P10, P50, P90)
    \item Tornado sensitivity chart ranking input parameters by NPV impact
\end{itemize}

% -----------------------------------------------------------
\subsection{Machine Learning Grade Prediction}
\label{sec:ml}

A supervised learning model predicts gold grade from spatial and geological features:

\begin{itemize}[itemsep=0.1em]
    \item \textbf{Features (11):} $x, y, z$ coordinates; density; RQD; $xy$, $xz$, $yz$ interactions; $r = \sqrt{x^2+y^2+z^2}$; depth ($|z|$); distance from centroid.
    \item \textbf{Models:} Random Forest (100 trees) and Gradient Boosting (100 estimators).
    \item \textbf{Uncertainty:} For Random Forest, per-block prediction variance is computed from individual tree predictions, providing a spatial uncertainty map.
\end{itemize}

\begin{equation}
    \hat{\sigma}^2(\mathbf{x}) = \frac{1}{T} \sum_{t=1}^{T} \left( \hat{g}_t(\mathbf{x}) - \bar{g}(\mathbf{x}) \right)^2
\end{equation}

where $\hat{g}_t(\mathbf{x})$ is the prediction of tree $t$ and $\bar{g}(\mathbf{x})$ is the ensemble mean.

% -----------------------------------------------------------
\subsection{Production Scheduling}
\label{sec:scheduling}

The production scheduler assigns optimized stopes to annual mining periods using a grade-priority strategy:

\begin{enumerate}
    \item Sort all selected stopes by average grade in descending order.
    \item Pack stopes into periods constrained by annual mining rate $M$ (default: 500{,}000\,t/year).
    \item For each period $t$, compute:
    \begin{itemize}
        \item Revenue: $R_t = \sum_{i \in P_t} m_i \times r \times p$
        \item Operating expenditure: $\text{OpEx}_t = \sum_{i \in P_t} T_i \times (C_m + C_p + C_{\text{ga}})$
        \item Development cost allocation: $D_t = C_{\text{dev}} / N_{\text{periods}}$
        \item Net cash flow: $\text{CF}_t = R_t - \text{OpEx}_t - D_t$
    \end{itemize}
    \item Apply Discounted Cash Flow analysis:
\end{enumerate}

\begin{equation}
    \text{DCF-NPV} = \sum_{t=1}^{N} \frac{\text{CF}_t}{(1 + d)^t}
\end{equation}

where $d$ is the discount rate (default 8\%) and $N$ is the mine life in years. The high-grade-first strategy maximizes early cash flows, improving the discounted NPV relative to alternative sequencing approaches.

% -----------------------------------------------------------
\subsection{Geotechnical Stability Assessment}
\label{sec:geotech}

\subsubsection{Stope Stability --- Mathews Method}

Each stope is assessed using a simplified Mathews stability graph approach \cite{mathews1981,potvin1988}:

\begin{enumerate}
    \item \textbf{In-situ stress:} $\sigma_v = \gamma \cdot z$ (vertical), $\sigma_h = K \cdot \sigma_v$ (horizontal, $K = 1.5$), where $\gamma = 0.027$\,MN/m$^3$.
    \item \textbf{Stability Number:} $N_s = \text{UCS} / \sigma_{\max}$, where UCS $\sim \mathcal{N}(120, 25^2)$\,MPa.
    \item \textbf{Risk classification:}
    \begin{itemize}
        \item $N_s > 10$: Stable (Low risk)
        \item $5 < N_s \leq 10$: Potentially Unstable (Medium risk)
        \item $2 < N_s \leq 5$: Unstable (High risk)
        \item $N_s \leq 2$: Caving Expected (Critical risk)
    \end{itemize}
    \item \textbf{Mathews $N$:} Calculated from Q-system rock mass quality, with stress, joint orientation, and gravity adjustment factors. Zone classification based on $N / HR$ ratio (Stable, Transitional, Caving).
\end{enumerate}

\subsubsection{Pillar Assessment}

Adjacent stope pairs are identified and pillar stability evaluated:

\begin{equation}
    \text{FoS} = \frac{\sigma_{\text{strength}}}{\sigma_{\text{stress}}} = \frac{0.6 \cdot \text{UCS} \cdot \sqrt{W_p / H_p}}{\gamma \cdot z \cdot (1 + W_s / W_p)}
\end{equation}

where $W_p$ is pillar width, $H_p$ is pillar height, and $W_s$ is stope width. Classification: Adequate (FoS $> 1.5$), Marginal ($1.0 < $ FoS $\leq 1.5$), Inadequate (FoS $\leq 1.0$).

% ============================================================
\section{System Architecture}
\label{sec:architecture}

\subsection{Technology Stack}

\begin{table}[H]
\centering
\caption{Technology stack}
\begin{tabular}{@{}ll@{}}
\toprule
\textbf{Component} & \textbf{Technology} \\
\midrule
Backend & Flask (Python 3.11+) \\
Frontend & HTML5, CSS3, JavaScript \\
Visualization & Plotly.js (interactive 3D) \\
Optimization & PuLP 4.x + HiGHS solver \\
Machine Learning & scikit-learn (Random Forest, GBM) \\
Numerical Computing & NumPy, Pandas, SciPy \\
Deployment & Render (cloud), Gunicorn (WSGI) \\
\bottomrule
\end{tabular}
\end{table}

\subsection{Web Application Design}

The application follows a RESTful API architecture:

\begin{itemize}[itemsep=0.2em]
    \item \textbf{Backend:} Flask serves 11 API endpoints, each accepting JSON parameters and returning computed results. An in-memory cache stores intermediate results (block model, optimization output) across sequential analysis steps.
    \item \textbf{Frontend:} A responsive, dark-themed single-page application with 9 interactive tabs. All charts are rendered client-side using Plotly.js for fast, GPU-accelerated visualization. Per-tab inline loading indicators provide non-blocking UX during long computations.
    \item \textbf{Landing page:} Features section, data calibration details, and an industry comparison table highlighting differences from traditional workflows.
\end{itemize}

\subsection{API Endpoints}

\begin{table}[H]
\centering
\caption{REST API endpoints}
\small
\begin{tabular}{@{}lll@{}}
\toprule
\textbf{Endpoint} & \textbf{Method} & \textbf{Description} \\
\midrule
\texttt{/} & GET & Landing page + dashboard \\
\texttt{/api/generate-block-model} & POST & Generate block model, evaluate economics \\
\texttt{/api/cross-section} & POST & Grade heatmap at given elevation \\
\texttt{/api/optimize} & POST & Floating stope + BILP optimization \\
\texttt{/api/mine-layout} & POST & Development network generation \\
\texttt{/api/economics} & POST & NPV, cost breakdown, sensitivity \\
\texttt{/api/monte-carlo} & POST & Monte Carlo NPV simulation \\
\texttt{/api/ml-predict} & POST & ML grade prediction model \\
\texttt{/api/ml-uncertainty} & POST & Prediction uncertainty at elevation \\
\texttt{/api/schedule} & POST & Production schedule with DCF NPV \\
\texttt{/api/geotechnical} & POST & Geotechnical stability assessment \\
\bottomrule
\end{tabular}
\end{table}

% ============================================================
\section{Results}
\label{sec:results}

\subsection{Block Model Statistics}

The generated block model exhibits the expected characteristics of an Archean lode gold deposit:

\begin{table}[H]
\centering
\caption{Block model summary statistics (default parameters)}
\begin{tabular}{@{}lr@{}}
\toprule
\textbf{Metric} & \textbf{Value} \\
\midrule
Total blocks & 24{,}000 \\
Mean grade & 0.94\,g/t \\
Median grade & 0.39\,g/t \\
Maximum grade & 39.4\,g/t \\
Coefficient of Variation & 181.5\% \\
Ore blocks ($> g_{\text{cutoff}}$) & $\sim$10{,}700 (44.9\%) \\
\bottomrule
\end{tabular}
\end{table}

The mean-to-median ratio of $\sim$2.4 and CV of 181.5\% confirm the strong positive skew characteristic of gold grade distributions, in contrast to the near-symmetric distributions typical of copper or iron ore deposits.

\subsection{Optimization Results}

With default stope dimension parameters:

\begin{table}[H]
\centering
\caption{Stope optimization results (default parameters)}
\begin{tabular}{@{}lr@{}}
\toprule
\textbf{Metric} & \textbf{Value} \\
\midrule
Candidate stopes generated & $\sim$3{,}400 \\
Optimal stopes selected & $\sim$420 \\
Total stope value & $\sim$\$380M \\
Total ore tonnage & $\sim$1.2M\,t \\
HiGHS solver time & $<$30\,s \\
Optimality gap & 0\% (provably optimal) \\
\bottomrule
\end{tabular}
\end{table}

\subsection{Comparison with Traditional Methods}

\begin{table}[H]
\centering
\caption{MineOptima vs traditional underground mine design practice}
\small
\begin{tabular}{@{}p{3cm}p{5cm}p{5cm}@{}}
\toprule
\textbf{Aspect} & \textbf{Traditional Practice} & \textbf{MineOptima} \\
\midrule
Stope layout & Manual trial-and-error in Surpac/Deswik (days--weeks) & BILP with provably optimal solution (seconds) \\
Grade estimation & Kriging/IDW with limited uncertainty & Random Forest with per-block uncertainty \\
Risk analysis & Deterministic single-scenario NPV & 1{,}000-iteration Monte Carlo with full distribution \\
Workflow & Disconnected tools requiring specialists & Integrated pipeline, single web interface \\
Visualization & Expensive CAD licenses & Browser-based, accessible to any stakeholder \\
Scheduling & Manual Gantt charts & Automated grade-priority with DCF \\
Geotechnical & Separate stability analysis & Integrated Mathews method + pillar FoS \\
\bottomrule
\end{tabular}
\end{table}

% ============================================================
\section{Project Structure}
\label{sec:structure}

\begin{lstlisting}[caption={Project directory structure}]
Automated_Underground_Mine_design/
|-- flask_app.py                  # Flask backend (REST API)
|-- requirements.txt              # Python dependencies
|-- README.md                     # Project documentation
|-- report.tex                    # This LaTeX report
|-- src/
|   |-- block_model_generator.py  # Synthetic gold deposit model
|   |-- economic_evaluator.py     # Block valuation & classification
|   |-- stope_generator.py        # Floating stope algorithm
|   |-- stope_optimizer.py        # BILP optimization (PuLP+HiGHS)
|   |-- mine_layout.py            # Development network generation
|   |-- monte_carlo.py            # Monte Carlo NPV simulation
|   |-- grade_prediction.py       # ML grade prediction (RF/GBM)
|   |-- production_scheduler.py   # Production scheduling + DCF
|   |-- geotechnical.py           # Geotechnical stability assessment
|-- templates/
|   |-- index.html                # HTML template (landing + dashboard)
|-- static/
|   |-- css/style.css             # Dark-themed responsive CSS
|   |-- js/app.js                 # Frontend logic, Plotly charts
|-- data/
    |-- block_model.csv           # Generated block model
    |-- economic_block_model.csv  # With economic values
\end{lstlisting}

% ============================================================
\section{Future Works}
\label{sec:future}

The following enhancements would further increase the platform's value for research and industry application:

\subsection{Multi-Scenario Comparison}
Allow users to save and compare multiple optimization runs side-by-side with different parameters (gold price, cutoff grade, stope dimensions). This would enable rapid what-if analysis for strategic mine planning, similar to the scenario management capabilities in commercial tools like Deswik.Stope.

\subsection{Real Data Integration (CSV Upload)}
Support importing real block model data from CSV or mining software exports (Surpac, Vulcan, Leapfrog). This would allow the platform to work with actual orebody data from exploration drilling programs, extending beyond synthetic models.

\subsection{Geostatistical Grade Estimation}
Integrate kriging methods (Ordinary Kriging, Simple Kriging, Indicator Kriging) alongside the existing ML-based grade prediction. Variogram modeling ($\gamma(h)$) and cross-validation would provide industry-standard geostatistical estimates compliant with JORC/NI~43-101 reporting standards.

\subsection{Backfill Optimization}
Model cemented rockfill (CRF) and paste fill requirements for extracted stopes. Optimize the fill sequence to maintain ground stability while maximizing ore recovery through secondary stoping---a critical consideration in long-hole open stoping operations.

\subsection{Ventilation Network Design}
Automatically generate a ventilation network layout based on the development and stope geometry. Calculate airflow requirements, fan sizing, and ensure compliance with regulatory limits (diesel particulate matter, heat stress) for underground personnel and equipment.

\subsection{Real-Time Production Tracking}
Add a production dashboard that tracks actual vs.\ planned tonnage, grade, and metal production per period. This would transform the tool from a planning platform into an operational monitoring system, enabling grade control reconciliation.

\subsection{3D Geological Modeling}
Integrate implicit geological modeling (radial basis functions or signed distance fields) to define ore zone boundaries, fault planes, and lithological contacts. This geological context would improve both stope generation constraints and geotechnical assessments.

\subsection{Multi-Objective Optimization}
Extend the BILP formulation to handle multiple objectives simultaneously using $\epsilon$-constraint or weighted-sum methods:

\begin{equation}
    \text{Maximize:} \quad \alpha_1 \cdot \text{NPV} - \alpha_2 \cdot \text{Risk} - \alpha_3 \cdot C_{\text{dev}}
\end{equation}

Pareto frontier visualization would help mine planners balance competing goals (e.g., maximize NPV while minimizing geotechnical risk or development cost).

\subsection{Cloud-Native Deployment}
Containerize the application with Docker, implement a job queue (Celery + Redis) for long-running optimizations, and add user authentication (OAuth2). This would support multi-user access, persistent sessions, and enterprise deployment scenarios.

% ============================================================
\section{Conclusion}
\label{sec:conclusion}

MineOptima demonstrates that modern computational methods---operations research, machine learning, and stochastic simulation---can be integrated into a single, accessible web platform for automated underground mine design. The system addresses the specific challenges of gold deposit modeling (nugget effect, log-normal grade distributions, narrow vein structures) while providing provably optimal stope layouts, probabilistic risk quantification, and integrated geotechnical assessment.

The platform reduces the mine design cycle from days or weeks of manual iteration to under 60 seconds of automated computation, while producing solutions that are mathematically guaranteed to be optimal (within the BILP formulation). By making these capabilities available through a browser-based interface, the system democratizes access to advanced mine planning tools that traditionally require expensive commercial software licenses and specialist expertise.

% ============================================================
\begin{thebibliography}{9}

\bibitem{journel1978}
A.G. Journel and Ch.J. Huijbregts, \textit{Mining Geostatistics}, Academic Press, 1978.

\bibitem{mathews1981}
K.E. Mathews, E. Hoek, D.C. Wyllie, and S.B.V. Stewart, ``Prediction of stable excavation spans for mining at depths below 1000 metres in hard rock,'' \textit{CANMET Report DSS Serial No.\ OSQ80-00081}, 1981.

\bibitem{potvin1988}
Y. Potvin, ``Empirical open stope design in Canada,'' PhD Thesis, University of British Columbia, 1988.

\bibitem{beasley1998}
J.E. Beasley, ``Heuristic algorithms for the unconstrained binary quadratic programming problem,'' Working Paper, Management School, Imperial College, London, 1998.

\bibitem{alford2007}
C. Alford, M. Brazil, and D.H. Lee, ``Optimisation in underground mining,'' in \textit{Handbook of Operations Research in Natural Resources}, Springer, 2007, pp.\ 561--577.

\bibitem{ataee2008}
M. Ataee-pour, ``A critical survey of the existing stope layout optimization techniques,'' \textit{Journal of Mining Science}, vol.\ 41, no.\ 5, pp.\ 447--461, 2005.

\bibitem{villalba2014}
E.A. Villalba Matamoros and M., ``Underground mine planning: stope layout optimisation under geological uncertainty using genetic algorithm,'' in \textit{Proc.\ Mine Planning and Equipment Selection}, 2014.

\bibitem{wgc2023}
World Gold Council, ``Gold mine cost benchmarks,'' \textit{All-In Sustaining Cost Report}, 2023. Available: \url{https://www.gold.org/}

\end{thebibliography}

\end{document}
